\documentclass[10pt,a4paper]{amsart}
\usepackage[margin=19mm]{geometry}
\usepackage{amsmath,amssymb,mathtools}
\usepackage[scr=rsfs,cal=euler]{mathalfa}
\usepackage{xfrac}
\usepackage[unicode,hidelinks,bookmarks=false]{hyperref}
\newcommand{\ie}{i.e.\ }
\newcommand{\cf}{cf.\ }
\newcommand{\resp}{respectively\ }
\newcommand{\Subsection}{\subsection}
\providecommand{\nobreakdash}{\nobreak\hbox{-}}
\setlength{\emergencystretch}{2em}
\multlinegap0pt
\usepackage{bm}
\usepackage{bbm}
\usepackage{dsfont}
\usepackage{xspace}
\usepackage{cancel}
\usepackage{mathtools}
\usepackage{amsthm}
\makeatletter
\@ifundefined{theorem}{\newtheorem{theorem}{Theorem}}{}
\@ifundefined{definition}{\newtheorem{definition}[theorem]{Definition}}{}
\@ifundefined{lemma}{\newtheorem{lemma}[theorem]{Lemma}}{}
\@ifundefined{proposition}{\newtheorem{proposition}[theorem]{Proposition}}{}
\makeatother
\newcommand{\N}{\mathbb{N}}
\newcommand{\Z}{{\mathbb Z}}
\newcommand{\akp}{\mathcal{A}_k(P)}
\newcommand{\bfx}{\mathbf{x}}
\newcommand{\bfy}{\mathbf{y}}
\newcommand{\cala}{\mathcal{A}}
\newcommand{\calc}{\mathcal{C}}
\newcommand{\inverse}{^{-1}}
\newcommand{\ra}{\rightarrow}
\title[Branching rules of minuscule representations via a new parti]{Branching rules of minuscule representations via a new partial order}
\author{R. M. Green, Tianyuan Xu}
\date{Source: arXiv:2402.06732v3 (CC BY 4.0)}
\begin{document}
\maketitle

\noindent\emph{Source and licence.} Branching rules of minuscule representations via a new partial order. Version arXiv:2402.06732v3 (CC BY 4.0), distributed under the Creative Commons Attribution 4.0 International licence, as recorded in the arXiv licence metadata for this exact version and shown on the abstract page https://arxiv.org/abs/2402.06732v3. Full rights notice: licenses/arxiv-2402.06732-CC-BY-4.0.txt.\par\smallskip

\noindent\textit{Editorial scope.} Counted unit: the Proposition whose source label is \texttt{prop:Dspin} in the article (source0.tex lines 982--987, source label \texttt{prop:Dspin}), delivered together with the article's own complete original proof of it (source0.tex lines 989--1059, one \texttt{proof} environment of 71 lines). No mathematics is written, repaired or re-derived: statement and proof are the article's own text line for line. Required context retained uncounted: prop:basics (lines 492--503); def:cak (lines 575--583); lem:cak (lines 598--616); prop:jab\_iso (lines 658--661); lem:jab\_antichain (lines 899--904). Every definition, display and label of those lines is retained unchanged. Omitted from this package and not redistributed: 1--491, 504--574, 584--597, 617--657, 662--898, 905--981, 988--988, 1060--2001. Nothing inside a retained excerpt is omitted or altered: every retained body is delivered exactly as the article's lines are written, so no omission receipt of any kind accompanies this package. Citations: the retained excerpts carry no citation command at all, so no bibliography entry of the article is transcribed and no reference is redistributed. Reference adaptation: none was needed and none was made; every reference inside the delivered unit resolves inside it, so no unresolved reference or citation is produced. Printed slips kept literal: no printed word, symbol, hypothesis or number of the counted statement or of its proof was altered. Layout and class: the article's own class is \texttt{amsart} with packages including babel, amsmath, graphicx, hyperref, amsthm, amssymb, fullpage, booktabs, tabularx, amscd, subfig, xspace, tikz; the wrapper is the lane's pinned \texttt{amsart} editorial wrapper at 10pt on A4, which restates the theorem-like environments the retained text needs and adds the article's own macro definitions that the retained text needs (\texttt{\textbackslash N}, \texttt{\textbackslash Z}, \texttt{\textbackslash akp}, \texttt{\textbackslash bfx}, \texttt{\textbackslash bfy}, \texttt{\textbackslash cala}, \texttt{\textbackslash calc}, \texttt{\textbackslash inverse}, \texttt{\textbackslash ra}), with extra wrapper packages (bm, bbm, dsfont, xspace, cancel, mathtools). The delivered numbering is the unit's own. Assets: no retained span contains a figure environment, a graphics-inclusion command, a PDF-page inclusion, a table or a TikZ picture; no image file is copied into this package, the package declares no asset dependency and no image file is redistributed with it. Licence: version arXiv:2402.06732v3 of this article is distributed under the Creative Commons Attribution 4.0 International licence, as recorded in the arXiv licence metadata for this exact version and shown on the abstract page https://arxiv.org/abs/2402.06732v3. A full rights notice is delivered at licenses/arxiv-2402.06732-CC-BY-4.0.txt. Characters: every retained excerpt is pure ASCII, so the delivered engine, XeLaTeX with its default Latin Modern text font, reports no missing character.
\begin{proposition}\label{prop:basics}
Let $A,B\in \akp$. 
\begin{itemize} 
\item [(i)]  If $A \leq_k B$, then the elements of $A = \{a_1, a_2, \cdots,
     a_k\}$ and $B = \{b_1, b_2, \cdots, b_k\}$ can be ordered in such a way
     that $a_i \leq_P b_i$ for all $1 \leq i \leq k$.
 \item [(ii)] The elements $A$ and $B$ are in a covering relation
      $A\lessdot_{\cala_k(P)} B$ if and only if we have both (1) $A \prec_k B$
      and (2) the unique elements $a \in A \backslash B$ and $b \in B \backslash
      A$ satisfy $a \lessdot_P b$.
\end{itemize}
\end{proposition}

\begin{definition}\label{def:cak}
For any $k,n\in \N$ such that $k \leq n$, we define $\calc(n,k)$ to be the poset
consisting of strictly increasing length-$k$ sequences with entries from $[n]$,
ordered by coordinatewise comparison: for sequences $\bfx = (x_1, x_2, \cdots,
x_k)$ and $\bfy = (y_1, y_2, \cdots, y_k)$ in $\calc(n, k)$, we define
$\bfx\le_{\calc(n,k)} \bfy$ if $x_i\le y_i$ for all $1\le i\le k$.  We
abbreviate the notation $\le_{\calc(n,k)}$ to $\le_\calc$ and
$\lessdot_{\calc(n,k)}$ to $\lessdot_\calc$.
\end{definition}

\begin{lemma}\label{lem:cak} 
Let $k,n$ be nonnegative integers such that $k \leq n$. 
\begin{itemize}
\item[{\rm (i)}] 
The function $\rho : \calc(n, k) \rightarrow \N$ defined by 
\[
    \rho( (x_1,\cdots,x_k)) = x_1 + x_2 + \cdots + x_k
\]
is a rank function on $\calc(n,k)$.  In other words, if $\bfx, \bfy \in \calc(n,
k)$ satisfy $\bfx \leq_\calc \bfy$, then we have $\rho(\bfx) \leq \rho(\bfy)$,
with $\rho(\bfy) = \rho(\bfx) + 1$ if and only if $\bfx \lessdot_{\calc} \bfy$.

\item[{\rm (ii)}]
Two elements $\bfx=(x_1,\cdots,x_k)$ and $\bfy=(y_1,\cdots, y_k)$  in $\calc(n,
k)$  are in a covering relation $\bfx\lessdot_\calc \bfy$ if and only if there
exists $i\in [k]$ such that $y_i=x_i+1$ and $y_j=x_j$ for all $j\in
[k]\setminus\{i\}$.
\end{itemize}
\end{lemma}

\begin{proposition}\label{prop:jab_iso}
For any $a,b\in \N$, the posets  $J([a] \times [b])$ and $\calc(a+b, b)$ are
isomorphic.\qed 
\end{proposition}

\begin{lemma}
\label{lem:jab_antichain}
Let $P=[a]\times[b]$. Then a subset of $P$ forms an antichain in $P$ if and
only if it can be written in the form $A=\{(x_1,y_1), \cdots, (x_k,y_k)\}$ where
$x_1 < \cdots < x_k$ and $y_1 > \cdots > y_k$.\qed
\end{lemma}

\begin{proposition}
    \label{prop:Dspin}
    If $n,k\in \N$ satisfy $0\le k\le \lfloor (n+2)/2\rfloor$, then we have
    isomorphisms of posets \[ \cala_k(J([n]\times[2]))\cong
    \cala_k(\calc(n+2,2))\cong \calc(n+2,2k).  \]
\end{proposition}

\begin{proof}
Since $J([n]\times [2])\cong \calc(n+2,2)$ by Proposition \ref{prop:jab_iso}, it
suffices to prove $\cala_k(\calc(n+2,2))\cong \calc(n+2,2k)$. Let
$P=\calc(n+2,2)$ and write every element of $P$ in the form $(x,y)$ where $x<y$
as in Definition \ref{def:cak}. Then we may view $P$ as embedded in the lattice
$\Z_+^2$ as we did the poset $[a]\times[b]$ in the previous two sections, because two elements $p_1=(x_1,y_1)$ and
$p_2=(x_2,y_2)$ in $P$ satisfy $p_1\le_\calc p_2$ if and only if $x_1\le x_2$
and $y_1\le y_2$. It follows that the characterization of antichains in
$[a]\times [b]$ given in Lemma \ref{lem:jab_antichain}
holds for the poset $P=\calc(n+2,2)$ as well, so that
every antichain of $P$ of size $k$ can be written uniquely as 
\[
    A = \{ (x_1, y_1), \ (x_2, y_2), \ \cdots, \ (x_k, y_k)\}
\] 
where $x_1 < x_2 < \cdots < x_k$ and $y_1 > y_2 > \cdots > y_k$.
By assumption we have $x_k<y_k$, so we have 
\[
x_1<x_2<\cdots<x_k<y_k<\cdots<y_2<y_1. 
\] 
It follows that there is a function $\phi : \cala_k(P) \ra \calc(n+2, 2k)$ given
by $$ \phi(A) = (x_1, x_2, \cdots, x_k, y_k, \cdots, y_2, y_1).$$  Furthermore,
$\phi$ is a bijection, because the assignment 
\[
    (x_1, x_2,\cdots, x_k, y_k, y_{k-1},\cdots, y_1) \mapsto \{(x_1,y_1), \cdots, (x_k,y_k)\}
\]
is clearly a two-sided inverse of $\phi$ that takes $\calc(n+2,2k)$ to $\akp$.

We claim that $\phi$ is an isomorphism of posets. To see this, it is enough to
prove that $\phi$ and $\phi\inverse$ preserve covering relations. Let $A \in
\cala_k(P)$, and write $A = \{ (x_1, y_1), \ (x_2, y_2), \ \cdots, \ (x_k,
y_k)\}$, where $x_1 < x_2 < \cdots < x_k$ and $y_1 > y_2 > \cdots > y_k$.

Suppose that we have $A' \lessdot_{\cala_k(P)} A$ for some $A'\in \akp$.
Proposition \ref{prop:basics} (ii) implies that $A'$ can be obtained from $A$ by
replacing one of the $(x_i, y_i)$ by $(x'_i, y'_i)$, where $(x'_i, y'_i)
\lessdot_P (x_i, y_i)$. By Lemma \ref{lem:cak} (ii), it follows that we
either have (a) $x_i'=x_i$ and $y_i'=y_i-1$, or (b) $x_i'=x_i-1$, $y_i'=y_i$.
In the first case, we have $x_1<\cdots< x_i'=x_i< \cdots<x_k$, so that we must
have 
\[ x_1<\cdots<x_i'<\cdots< x_k<y_k<\cdots <y_i'<\cdots<y_1 \] 
by the argument from the first paragraph of this proof. It follows that 
\begin{eqnarray*} 
\varphi(A')&=&(x_1,\cdots,x_i, \cdots,x_k, y_k, \cdots,y_i', \cdots,  y_1) \\ 
      &\lessdot&(x_1,\cdots,x_i, \cdots,x_k, y_k, \cdots,y_i, \cdots, y_1)\\ 
      &=&\varphi(A) \end{eqnarray*} 
in $\calc(n+2,2k)$. Moreover, since $y_i'=y_i-1$, the relation
$\varphi(A')<\varphi(A)$ is a covering relation in $\calc(n+2,2k)$ by Lemma
\ref{lem:cak} (ii).  Similarly, a symmetric argument shows that in the second
case we have 
\begin{eqnarray*} 
\varphi(A')&=&(x_1,\cdots,x_i', \cdots,x_k,y_k, \cdots,y_i, \cdots,  y_1) \\ 
&\lessdot&(x_1,\cdots,x_i, \cdots,x_k, y_k, \cdots,y_i, \cdots,  y_1)\\ 
&=&\varphi(A) 
\end{eqnarray*} 
in $\calc(n+2,2k)$. 

Conversely, suppose that $\varphi(A')\lessdot_\calc \varphi(A)$ for some $A'\in
\akp$.  Lemma \ref{lem:cak} (ii) implies that one of the following conditions
must hold for some $i\in [k]$:

\begin{enumerate}
\item we have $\varphi(A')=(x_1, \cdots, x_i'=x_i-1, \cdots, x_k,y_k,\cdots, y_1)$; 
\item we have $\varphi(A)=(x_1, \cdots, x_k,y_k,\cdots,y_i'=y_i-1,\cdots, y_1)$. 
\end{enumerate} 

In the first case, we can use the inverse of $\varphi$ mentioned earlier to
recover $A'$ as the set $$ A'=\{(x_j,y_j):j\in [k], j\neq i\}\cup\{(x'_i,y_i)\}
,$$ where $(x'_i,y_i)\lessdot_P (x_i,y_i)$ by Lemma \ref{lem:cak} (ii); therefore $A'\lessdot_{\akp} A$ by Proposition \ref{prop:basics} (ii).  A
similar argument shows that $A'\lessdot_{\akp} A$ in the second case, and we are
done.
 \end{proof}

\end{document}
